\documentclass[11pt]{article}
\usepackage[margin=1in]{geometry}
\usepackage{amsmath,amssymb,amsthm}
\usepackage{graphicx}
\usepackage{hyperref}
\hypersetup{hidelinks}

\newtheorem{theorem}{Theorem}
\newtheorem{remark}{Remark}

\newcommand{\R}{\mathbb R}
\newcommand{\E}{\mathbb E}
\newcommand{\Pp}{\mathbb P}
\newcommand{\F}{\mathcal F}
\newcommand{\G}{\mathcal G}
\newcommand{\calL}{\mathcal L}
\newcommand{\gam}{\gamma}

\title{A Gaussian subcase of the Veraar--Weis Theorem 4.2 tail question}
\author{arXiv automatic math run \texttt{fa\_banach\_001}}
\date{2026-06-18}

\begin{document}
\maketitle

\section*{Source and target}

The source paper is Veraar and Weis,
\emph{A note on maximal estimates for stochastic convolutions},
arXiv:1004.5061 \cite{VW1004}.  The paper proves maximal estimates for
stochastic convolutions
\[
        S\diamond G(t)=\int_0^t S(t-s)G(s)\,dW_H(s)
\]
under an \(H^\infty\)-calculus assumption on the generator.  The first
exponential-tail question in the paper is answered immediately by
Theorem~2.4 for UMD spaces of type \(2\).  The surviving open target is the
paragraph after Theorem~4.2:
\begin{quote}
At this moment we do not know if there are exponential tail estimates in the
general setting of Theorem~4.2.  However, also in this setting there is some
hope that \(C_{p,Y}^{\gamma}\leq C_X\sqrt p\) for \(p\) large, and by the
argument in Theorem~2.4 this would yield exponential tail estimates again.
\end{quote}
The same paragraph records that Cox--Veraar obtained only
\(C_{p,Y}^{\gamma}\lesssim_X p\), which gives exponential estimates but not
the expected quadratic exponential estimates.

\begin{center}
\includegraphics[width=0.92\linewidth]{figures/open_problem_crop.png}
\end{center}

\section*{Result}

The full adapted problem remains open, but the quadratic tail holds in the
deterministic and initially independent subcase of Theorem~4.2.

\begin{theorem}
\label{thm:main}
Let \(X\) be a UMD Banach space with property \((\alpha)\), and let \(A\)
satisfy hypothesis \((H)\) of \cite{VW1004}.  Let \(S(t)=e^{tA}\).  Let
\((W_H(t))_{t\ge0}\) be an \(H\)-cylindrical Brownian motion, and suppose that
the filtration has the form
\[
        \F_t=\G\vee\sigma(W_H(s):0\le s\le t),
\]
where \(\G\) is independent of \(W_H\).  Let
\[
        G:\Omega\to \gam(\R_+;H,X)
\]
be \(\G\)-measurable and suppose that for some \(M>0\),
\[
        \|G(\omega)\|_{\gam(\R_+;H,X)}\le \sqrt M
        \quad\text{for a.e. }\omega .
\]
Then \(S\diamond G\) has the continuous modification supplied by
Theorem~4.2, and if \(K_2=K_2(A,X)\) denotes the \(p=2\) maximal-estimate
constant in Theorem~4.2, then for every \(u\ge0\),
\[
  \Pp\!\left(\sup_{t\ge0}\|S\diamond G(t)\|
       \ge K_2\sqrt M + u\right)
  \le \exp\!\left(-\frac{u^2}{2K_2^2M}\right).
\]
Consequently, for every \(\lambda\ge 2K_2\sqrt M\),
\[
  \Pp\!\left(\sup_{t\ge0}\|S\diamond G(t)\|\ge\lambda\right)
  \le \exp\!\left(-\frac{\lambda^2}{8K_2^2M}\right).
\]
\end{theorem}

\section*{Proof Intuition}

The obstruction in the adapted case is that the integrand can react to the
Brownian path; moment estimates then pass through Banach-space stochastic
integral constants.  If the integrand is known at time zero and independent
of the Brownian motion, the obstruction disappears.  Conditioning on \(\G\)
freezes \(G\), and the stochastic convolution becomes a centered Gaussian
process indexed by \(t\ge0\).

Theorem~4.2 at \(p=2\) gives an \(L^2\) bound for the supremum of that process.
For a centered Gaussian process, this same bound controls both the mean scale
and the weak variance parameter.  Borell's inequality then upgrades the
\(L^2\) maximal estimate to a quadratic tail.

\section*{Proof}

Fix a deterministic \(g\in\gam(\R_+;H,X)\).  Theorem~4.2 of \cite{VW1004},
with \(p=2\), gives the continuous modification and the estimate
\begin{equation}
\label{eq:p2}
 \left(\E_W\sup_{t\ge0}\|S\diamond g(t)\|^2\right)^{1/2}
 \le K_2\|g\|_{\gam(\R_+;H,X)} ,
\end{equation}
where \(\E_W\) denotes expectation over the Brownian motion and \(K_2\)
depends only on the constants in Theorem~4.2 for \(A\) and \(X\).

For deterministic \(g\), the process \(Z_t=S\diamond g(t)\) is centered
Gaussian.  Since the modification is continuous, the random variable
\[
        Z^*=\sup_{t\ge0}\|Z_t\|
\]
is the supremum of a separable centered Gaussian process after restricting
time to the rational points and using a countable norming set for the
separable closed span of the path.  Its weak variance parameter satisfies
\[
  \sigma_g^2
  =\sup_{t\ge0,\ \|x^*\|\le1}\E_W|x^*Z_t|^2
  \le \E_W (Z^*)^2
  \le K_2^2\|g\|_{\gam(\R_+;H,X)}^2
\]
by \eqref{eq:p2}.  Also
\[
        \E_W Z^*\le \left(\E_W (Z^*)^2\right)^{1/2}
        \le K_2\|g\|_{\gam(\R_+;H,X)} .
\]
Borell's inequality for separable centered Gaussian processes therefore gives
\[
  \Pp_W\!\left(
      Z^*\ge K_2\|g\|_{\gam(\R_+;H,X)}+u
    \right)
  \le
  \exp\!\left(
      -\frac{u^2}{2K_2^2\|g\|_{\gam(\R_+;H,X)}^2}
    \right)
\]
for \(u\ge0\), with the evident trivial interpretation when \(g=0\).

Now return to the random initially independent integrand \(G\).  Conditional
on \(\G\), the preceding deterministic argument applies to \(g=G(\omega)\).
Using \(\|G(\omega)\|_{\gam(\R_+;H,X)}\le\sqrt M\) gives, a.s.,
\[
  \Pp\!\left(
      \sup_{t\ge0}\|S\diamond G(t)\|\ge K_2\sqrt M+u
      \,\middle|\,\G
    \right)
  \le
  \exp\!\left(-\frac{u^2}{2K_2^2M}\right).
\]
Taking expectation in \(\G\) proves the first displayed estimate.  If
\(\lambda\ge2K_2\sqrt M\), put \(u=\lambda-K_2\sqrt M\); then
\(u\ge\lambda/2\), which gives the second estimate.

\section*{Scope and verification notes}

This packet gives a solved subcase, not a full answer to the source paper's
general adapted question.  The proof uses Gaussianity after conditioning on
the initial sigma-field.  For a genuinely adapted integrand there is no such
conditional Gaussian process, and the source paper's proposed moment route
still requires the open \(\sqrt p\)-growth of the relevant UMD
stochastic-integral constants.

The boundary is consistent with later literature.  Cox and Veraar
\cite[p.~28]{CV1107} prove \(p\)-independence phenomena for decoupling, but
explicitly state that the optimal \(\kappa_{p,X}/\sqrt p\) boundedness is open
already for \(X=L^q\), \(q\ne2\).  Later general UMD Burkholder--Davis--Gundy
theorems, such as Yaroslavtsev \cite{Y1807}, provide qualitative
Banach-valued BDG estimates but do not remove this \(\sqrt p\) constant
obstruction for the adapted Theorem~4.2 tail problem.

Local cheap-index search found no existing result packet for arXiv:1004.5061.
The exact source paragraph, \(C_{p,Y}^{\gamma}\), Cox--Veraar arXiv:1107.2218,
and Yaroslavtsev arXiv:1807.05573 were checked as bounded literature/status
search terms.  No full adapted quadratic-tail answer was found in that check.

\section*{Human-review recommendation}

Review the separability step in the application of Borell's inequality and the
conditional-use formulation for initially independent integrands.  The
mathematical dependence on the source theorem is only the \(p=2\) maximal
estimate from Theorem~4.2, so the result should be robust once those standard
Gaussian-process details are accepted.

\begin{thebibliography}{9}

\bibitem{VW1004}
M. Veraar and L. Weis.
\newblock \emph{A note on maximal estimates for stochastic convolutions}.
\newblock arXiv:1004.5061, 2010.

\bibitem{CV1107}
S. Cox and M. Veraar.
\newblock \emph{Vector-valued decoupling and the Burkholder--Davis--Gundy
inequality}.
\newblock arXiv:1107.2218, 2011.

\bibitem{Y1807}
I. Yaroslavtsev.
\newblock \emph{Burkholder--Davis--Gundy inequalities in UMD Banach spaces}.
\newblock arXiv:1807.05573, 2018.

\bibitem{Ledoux}
M. Ledoux.
\newblock \emph{The concentration of measure phenomenon}.
\newblock American Mathematical Society, 2001.

\end{thebibliography}

\end{document}
